\documentclass[10pt,a4paper]{amsart}
\usepackage[margin=19mm]{geometry}
\usepackage{amsmath,amssymb,mathtools}
\usepackage[scr=rsfs,cal=euler]{mathalfa}
\usepackage{xfrac}
\usepackage[unicode,hidelinks,bookmarks=false]{hyperref}
\newcommand{\ie}{i.e.\ }
\newcommand{\cf}{cf.\ }
\newcommand{\resp}{respectively\ }
\newcommand{\Subsection}{\subsection}
\providecommand{\nobreakdash}{\nobreak\hbox{-}}
\setlength{\emergencystretch}{2em}
\multlinegap0pt
\usepackage{bm}
\usepackage{bbm}
\usepackage{dsfont}
\usepackage{xspace}
\usepackage{cancel}
\usepackage{mathtools}
\usepackage{amsthm}
\makeatletter
\@ifundefined{observation}{\newtheorem{observation}{Observation}}{}
\@ifundefined{simplebox}{\newtheorem{simplebox}{Simplebox}}{}
\@ifundefined{theorem}{\newtheorem{theorem}{Theorem}}{}
\@ifundefined{claim}{\newtheorem{claim}[theorem]{Claim}}{}
\@ifundefined{lemma}{\newtheorem{lemma}[theorem]{Lemma}}{}
\makeatother
\DeclareMathOperator{\interior}{int}
\newcommand{\N}{\mathbb{N}}
\newcommand{\Q}{\mathbb{Q}}
\title[Normality in the square of the Sorgenfrey Line]{Normality in the square of the Sorgenfrey Line}
\author{Paul Szeptycki, Hongwei Wen}
\date{Source: arXiv:2511.12327v2 (CC BY 4.0)}
\begin{document}
\maketitle

\noindent\emph{Source and licence.} Normality in the square of the Sorgenfrey Line. Version arXiv:2511.12327v2 (CC BY 4.0), distributed under the Creative Commons Attribution 4.0 International licence, as recorded in the arXiv licence metadata for this exact version and shown on the abstract page https://arxiv.org/abs/2511.12327v2. Full rights notice: licenses/arxiv-2511.12327-CC-BY-4.0.txt.\par\smallskip

\noindent\textit{Editorial scope.} Counted unit: the Theorem whose source label is \texttt{None} in the article (source0.tex lines 484--486, source label \texttt{None}), delivered together with the article's own complete original proof of it (source0.tex lines 487--564, one \texttt{proof} environment of 78 lines). No mathematics is written, repaired or re-derived: statement and proof are the article's own text line for line. Required context retained uncounted: setE (lines 414--423). Every definition, display and label of those lines is retained unchanged. Omitted from this package and not redistributed: 1--413, 424--483, 565--695. Nothing inside a retained excerpt is omitted or altered: every retained body is delivered exactly as the article's lines are written, so no omission receipt of any kind accompanies this package. Citations: the retained excerpts carry no citation command at all, so no bibliography entry of the article is transcribed and no reference is redistributed. Reference adaptation: none was needed and none was made; every reference inside the delivered unit resolves inside it, so no unresolved reference or citation is produced. Printed slips kept literal: no printed word, symbol, hypothesis or number of the counted statement or of its proof was altered. Layout and class: the article's own class is \texttt{amsart} with packages including amssymb, amsmath, amsthm, babel, inputenc, fontenc, enumerate, mathtools, graphicx, subfigure, float, adjustbox, tikz, multirow; the wrapper is the lane's pinned \texttt{amsart} editorial wrapper at 10pt on A4, which restates the theorem-like environments the retained text needs and adds the article's own macro definitions that the retained text needs (\texttt{\textbackslash N}, \texttt{\textbackslash Q}, \texttt{\textbackslash interior}), with extra wrapper packages (bm, bbm, dsfont, xspace, cancel, mathtools). The delivered numbering is the unit's own. Assets: no retained span contains a figure environment, a graphics-inclusion command, a PDF-page inclusion, a table or a TikZ picture; no image file is copied into this package, the package declares no asset dependency and no image file is redistributed with it. Licence: version arXiv:2511.12327v2 of this article is distributed under the Creative Commons Attribution 4.0 International licence, as recorded in the arXiv licence metadata for this exact version and shown on the abstract page https://arxiv.org/abs/2511.12327v2. A full rights notice is delivered at licenses/arxiv-2511.12327-CC-BY-4.0.txt. Characters: every retained excerpt is pure ASCII, so the delivered engine, XeLaTeX with its default Latin Modern text font, reports no missing character.
\begin{theorem}\label{setE}
    Under CH, there exists a set $E$ such that 
    \begin{enumerate}
        \item For any continuous function $f$ from a $G_\delta$ subset of ${\mathbb R}^n$ into $E$, $(n<\omega)$ there is an $\alpha<2^\omega$ such that
        \[
        \forall \beta\geq \alpha, f(\{r_\gamma:\gamma<\beta\}^n) \cap E \subseteq \{r_\gamma:\gamma<\beta\}
        \]
        \item $E$ contains $\Q$ and is concentrated on ${\mathbb Q}$. 
    \end{enumerate}
\end{theorem}

\begin{theorem}
    Let $E$ be any subset of ${\mathbb R}$ that admits no uncountable strictly decreasing functions (e.g., the set constructed in Theorem \ref{setE}), then $(E[\leq])^2$ is normal.
\end{theorem}

\begin{proof}
    Let $A,B$ be Sorgenfrey closed in $(E[\leq])^2$. Recall the shoelace lemma: to separate $A$ and $B$ it suffices to cover each set by a countable family of open sets whose closures are disjoint from the other closed set. By symmetry, it suffices to find a countable cover $\{U_n\}_{n<\omega}$ of $A$ such that $\overline{U_n}^\sigma\cap B=\emptyset$.  
    
    
   Recall we defined 
        $A^\prime\subseteq A$ to have the \textbf{covering property} (with respect to $B$) if there is a countable open cover $\{U_n\}$ such that $\overline{U_n}^\sigma\cap B=\emptyset$ 

    Clearly, we have: 
    \begin{lemma}\label{lem:lindelof}
        If $A^\prime\subseteq A$ is Lindel\"of, then $A^\prime$ has the covering property.
    \end{lemma}
    %\begin{proof}
      %  For any $a\in A$, there exists an $r_a>0$ such that the basic nbhd $U(a,r_a)$ of $a$ such that $U(a,r_a)\cap B=\emptyset$. Since $U(a,r_a)$ is also closed, it follows that $\overline{U(a,r_a)}\cap B=\emptyset$. Then since $A^\prime$ is Lindelof, the cover $\{U(a,r_a)\}_{a\in A^\prime}$ will be reduced to a countable cover, hence $A^\prime$ has the covering property.
   % \end{proof}
    For all $z\in B^c$ choose $n_z\in \N$ such that $U(z,1/n_z)\cap B=\emptyset$, this is possible since $B$ is Sorgenfrey closed.
 By letting $A_1= A\cap \overline{B}^E$ and $A_2= A\setminus \overline{B}^E$, we have 
    \begin{simplebox}
    \begin{observation}
    $A=A_1\cup A_2$
    \end{observation}
    \end{simplebox}
    \begin{claim}\label{clm:eu_open}
        $A_2$ has the covering property.
    \end{claim}
    \begin{proof}
        Note that $A_2 \subseteq E\setminus cl_E(B)=U$ and the latter is Euclidean open. Since $E^2$ is hereditarily Lindel\"of with respect to the Euclidean topology, $U$ can be written as the countable union of basic open sets $B_n$ such that $\overline{B_n}^\sigma\cap B=\emptyset$ for each $n$ so, since the Sorgenfrey topology is finer than the Euclidean, $A_2$ has the covering property. 
        %For every point $z\in E\setminus cl_E(B)$, There exists $w_z$ and $r_z>0$ such that $U(w_z,r_z)\subseteq E\setminus cl_E(B)$ and $z\in U(w_z,r_z)$. Hence 
        %$$\{\interior_E U(w_z,r_z) : z\in E\setminus cl_E(B)\}$$ will be reduced to a countable cover $\{\interior_E U(w_n,r_n)\}_{n\geq1}$ since Euclidean topology is hereditarily Lindelof. Then we have $\{U(w_n,r_n)\}_{n\geq 1}$ covers $E\setminus cl_E(B)\supseteq A_2$, and each $U(w_n,r_n)$ is also closed hence $\overline{U(w_n,r_n)}\cap B= U(w_n,r_n)\cap B=\emptyset$, $A_2$ has the covering property.
    \end{proof}
    \begin{claim}\label{clm:A1_cov}
        $A_1$ has the covering property.
    \end{claim}
    \begin{proof}  We will break down $A_1$ into countably many pieces, each of which is the graph of a monotone decreasing function and prove every piece has the covering property. Define
    \[A_q=\{ z=(x,y)\in A_1:x+y<q<x+y+1/n_z \}\]
    Denote $l_q= \{(x,y)\in X^2: x+y=q\}$, $A_q$ is just the collection of points $z$ which has the straight line $l_q$ passing through the Euclidean interior of its nbhd $U(z,1/n_z)$. We see that 
    \begin{simplebox}
        \begin{observation}
            $A_1= \cup_{q\in \Q}A_q$
        \end{observation}
    \end{simplebox} So now it suffices to show that each $A_q$ has the covering property. First note that 
    \begin{claim}(non-increasing property)\label{clm:diagonal}
    For all $z_1=(x_1,y_1),z_2=(x_2,y_2)\in A_q$, if $x_1< x_2$, then $y_2\leq y_1$. 
    %(contrapositive will also be used)

    \end{claim}
    \begin{proof}
        For otherwise we have $x_1<x_2$ and $y_1<y_2$. Since the line $l_q$ passes through both nbhds $U(z_1,1/n_{z_1}),U(z_2,1/n_{z_2})$, this means $z_2$ is in the Euclidean interior of $U(z_1,1/n_{z_1})$. However, $\interior_E U(z_1,1/n_{z_1})$ is a Euclidean open set, and since $z_2\in \overline{B}^{E}$, we have that $\interior_E U(z_1,1/n_{z_1})$ also contains a point from $B$, which is a contradiction, because we chose $n_{z_1}$ so that its nbhd $U(z_1,1/n_{z_1})$ is disjoint from $B$.
    \end{proof}
    \begin{claim}\label{clm:ccc}
        The sets $M_1=\{x\in p_0(A_q): |p_0^{-1}(x)|>1\}, M_2=\{x\in p_1(A_q):|p_1^{-1}(x)|>1\}$ are countable.
    \end{claim}
    \begin{proof}
         For each $x\in M_1$, fix $\alpha_x<\beta_x\in p_0^{-1}(x)$. And for any $x_1< x_2\in M_1$, we have $\beta_{x_2}\leq \alpha_{x_1}$ by the previous claim. And so the intervals $(\alpha_{x_1},\beta_{x_1})$ and $ (\alpha_{x_2},\beta_{x_2})$ are disjoint.  Hence, we have an uncountable set of pairwise disjoint open intervals on ${\mathbb R}$, contradiction.  The proof for $M_2$ is the same. 
    \end{proof}
     Let $L_1=p_0^{-1}(M_1),L_2=p_1^{-1}(M_2)$, and $L=L_1\cup L_2$, $K= A_q\setminus L$, we see that 
    \begin{simplebox}
        \begin{observation}
            $A_q= K\cup L$
        \end{observation}
    \end{simplebox}
    \begin{claim}
        $L$ has the covering property.
    \end{claim}
    \begin{proof}
        We know $L_1,L_2$ are Lindel\"of since by \ref{clm:ccc} each is a countable union of Lindel\"of subspaces. And by \ref{lem:lindelof}, $L$ has the covering property.
    \end{proof}
    Now it suffices to show $K$ has the covering property.
    \begin{claim}\label{clm:final}
        $K$ has the covering property
    \end{claim}
    \begin{proof}
        Note that $K$ is now a graph of some strictly decreasing function $g$. By the definition of $E$, such $g$ can only be defined on a countable set, hence the graph is countable hence Lindel\"of.
    \end{proof}
    This shows that each $A_q$ has the covering property completing the proof that $A_1$ has the covering property.
    \end{proof}
    
    Now with all the observations and claims, we obtain a countable cover of $A$ witnessing that $A$ has the covering property and so by the shoelace lemma, $A$ and $B$ can be separated by open sets, hence $E^2$ is normal.
\end{proof}

\end{document}
